% Generated draft: requires mathematical and source review.
\documentclass[11pt]{article}
\usepackage[margin=25mm]{geometry}
\usepackage{fontspec}
\setmainfont{Latin Modern Roman}
\usepackage{amsmath,amssymb,mathtools,unicode-math}
\setmathfont{Latin Modern Math}
\usepackage{enumitem,needspace,xurl,hyperref}
\setlength{\parindent}{0pt}
\setlength{\parskip}{5pt}
\setlength{\emergencystretch}{3em}
\newcommand{\N}{\mathbb{N}}
\newcommand{\Z}{\mathbb{Z}}
\newcommand{\Q}{\mathbb{Q}}
\newcommand{\R}{\mathbb{R}}
\newcommand{\C}{\mathbb{C}}
\newcommand{\F}{\mathbb{F}}
\newcommand{\A}{\mathbb{A}}
\newcommand{\PP}{\mathbb{P}}
\newcommand{\E}{\mathbb{E}}
\newcommand{\eps}{\varepsilon}
\newcommand{\dd}{\,\mathrm{d}}
\begin{document}
\section*{UnsolvedMath ID: 30006143}
\textbf{Source problem code:} catalogue entry
\subsection{Definitions and mathematical statement}
For an integer \(r\geq 1\), a line in \(\R^r\) means a one-dimensional linear subspace. A finite set \(\mathcal L\) of distinct lines is equiangular if there exists \(c\in[0,1)\) such that, for unit vectors \(u_L\in L\) chosen for each \(L\in\mathcal L\),
\[
\lvert\langle u_L,u_{L'}\rangle\rvert=c
\qquad\text{whenever }L,L'\in\mathcal L,\ L\ne L'.
\]
The absolute value makes this condition independent of the choices of signs of the unit vectors. The common value \(c\) may depend on the configuration; it is not prescribed in advance. Define
\[
N(r)=\max\{\lvert\mathcal L\rvert:\mathcal L\text{ is an equiangular set of lines in }\R^r\}.
\]
The supplied source states Gerzon's absolute bound
\[
N(r)\leq\binom{r+1}{2}=\frac{r(r+1)}2.
\]
The problem is to determine whether
\[
\bigl\{r\in\Z:r\geq1,\ N(r)=\tfrac{r(r+1)}2\bigr\}
\quad\text{is infinite}.
\]
Equivalently, does every integer \(R\geq1\) admit an integer \(r\geq R\) with an equiangular set of exactly \(r(r+1)/2\) lines in \(\R^r\)?

The original excerpt additionally describes the bound as tight up to a multiplicative constant; this is supplied background, distinct from the exact equality question. The supplied dated literature assessment reports equality in dimensions \(2,3,7,23\), nonattainment in infinitely many dimensions, and that the infinitude of equality dimensions remains open. These are supplied curation claims, not independently verified conclusions of this draft.
\subsection{Short English statement}
Are there arbitrarily large positive integer dimensions \(r\) in which one can find exactly \(r(r+1)/2\) distinct lines through the origin in \(\R^r\), all making the same angle with one another? The common angle may vary between configurations. The registry labels the problem partially solved; this draft preserves that supplied assessment without independently verifying it.
\subsection{Sources}
\begin{itemize}
\item Local UnsolvedMath registry, \texttt{data/unsolvedmath/problems.json}, permanent ID \(30006143\), source problem code \texttt{catalogue entry}, title ``Equality Cases for Gerzon's Equiangular-Line Bound.'' The supplied record labels the statement \texttt{corrected\_verified} and its status \texttt{partially\_solved}; its verification note says that only formula delimiters, notation, grammar, and typography were repaired, preserving mathematical content and scope. These are supplied registry annotations.
\item Source report: \emph{Graph Theory} (2025), Oberwolfach Reports, pp.\ 5--74, DOI \texttt{10.4171/owr/2025/1}. \url{https://doi.org/10.4171/owr/2025/1}. The registry describes these reports as documenting mathematical workshops and their open-problem discussions.
\item The supplied original excerpt associates the discussion with ``Equiangular lines via improved eigenvalue multiplicity,'' Matija Buci\'{c}, joint work with Igor Balla. The registry gives proposed year 2025 but does not identify a proposer. Its excerpt cites a reference labeled [6], whose bibliographic entry is not supplied.
\item Supplied literature source: \emph{Equiangular lines in Euclidean spaces}. \url{https://doi.org/10.1016/j.jcta.2015.09.008}.
\item Supplied literature source: \emph{Equiangular lines and Gerzon's bound}. \url{https://doi.org/10.1007/s10107-021-01638-x}.
\item Literature assessment checked at \texttt{2026-08-22}: this is supplied dated curation derived from the OWR record, not independently verified. It reports that Gerzon's bound is known, equality is attained in dimensions \(2,3,7,23\), infinitely many dimensions cannot attain equality, and whether equality occurs in infinitely many dimensions remains open. The supplied curation explicitly calls for independent verification of time-sensitive status and cited claims.
\end{itemize}
\end{document}
